%\begin{lemma} Машину $M$ с ленточным алфавитом $\Gamma$ и с 1-ой бесконечной в обе стороны лентой можно смоделировать на машине $M'$ с 2-мя бесконечными вправо лентами без замедления.
%\end{lemma}

%\begin{proof} "Перегнём" ленту пополам в произвольном месте (то есть получим две ленты бесконечные вправо, но ещё и соединённые между собой) и рассмотрим её как ленту с символами, представляющими собой пары символов. Тогда новый ленточный алфавит будет $\Gamma^2$. При работе с такой машиной будем запоминать, с какой стороны перегнутой ленты мы находимся, при помощи состояния. Если сверху — игнорируем второй символ в ячейке, если снизу — игнорируем первый символ.

%\noindent (Можно чуть иначе: хранить на ленте символ из пары <символ $\in \Gamma$, номер ленты $\in \{1, 2\}$>. Тогда новый ленточный алфавит будет представлять из себя $\Gamma \times \{1, 2\}$)
%\newline Замедления не будет, так как на каждой ленте есть головка.
%\end{proof}

\begin{lemma} Машину Тьюринга $M = \{\Sigma, \Gamma, Q, q_{\text{start}}, q_{\text{accept}}, q_{\text{reject}}, \delta\}$ с одной бесконечной в обе стороны лентой можно смоделировать на машине $M'$ с одной бесконечной вправо лентой без замедления.
\end{lemma}

\begin{proof} Рассмотрим ленту, ограниченную слева, с ленточным алфавитом $\Gamma^2$. Для моделирования одноленточной неограниченнной машины Тьюринга  раздвоим все её состояния $Q_{\text{new}} = Q \cup \text{copy}(Q)$: одни будут изменять только первые элементы пары из $\Gamma^2$, а их копии - только вторые. При этом изменим функцию перехода $\delta_{\text{new}}$, так чтобы инвертировать движение вправо/влево на состояниях из $\text{copy}(Q)$. Исправим в ней так же момент, когда головка находится в самой левой ячейке и хочет пойти влево: тогда ей нужно стоять на месте и перейти в состояние противоположного типа. Таким образом, $M' = \{\Sigma, \Gamma^2, Q_{\text{new}}, q_{\text{start}}, q_{\text{accept}}, q_{\text{reject}}, \delta_{\text{new}}\}$.

Это позволяет нам якобы сложить пополам неограниченную ленту и тем самым моделировать её лишь на одной ограниченной слева ленте. Поскольку каждому шагу $M$ соответствует один шаг на построенной нами $M'$ и наоборот, то замедление отсутствует. \end{proof}


%\begin{corollary} Машину $M$ с ленточным алфавитом $\Gamma$ и с $k$ бесконечными в обе стороны лентами можно смоделировать на машине $M'$ с $2k$ бесконечными вправо лентами без замедления.
%\end{corollary}

\begin{lemma} Машину $M$ с ленточным алфавитом $\Gamma$ и с $k$ бесконечными вправо лентами можно смоделировать на машине $M'$ с 1-ой бесконечной вправо лентой с замедлением в квадратичное число раз: $O(T(n)) \rightarrow O((T(n)^2)$, где $O(T(n))$ — время работы $M$ (верняя граница).
\end{lemma}

\begin{proof} \par
Пронумеруем на каждой $n\text{-ой}$ ленте каждую ячейку $i$ числом $n + i \cdot k$. Таким образом, чтобы понять номер исходной ленты по новому номеру ячейки, достаточно сделать $mod\ k$.

С учётом ограничения времени работы $M$ и, как следствие, ячеек, до которых можно физически дойти, будет конечное число — $O(T(n))$. Теперь это конечное число ячеек выложим номерами по возрастанию на новой одноленточной машине $M'$ с 1-ой бесконечной вправо лентой.

Ленточный алфавит на новой $M'$ будет $\Gamma \times \{0, 1\}$, где второе число значит, была ли головка на этой ячейке на её исходной ленте. Так как по номеру ячейки можно узнать номер её исходной ленты на $M$, значит, можно узнать, есть ли на этой ячейке головка.

%Теперь оценим замедление. Чтобы на одной ленте сместить головку на любой ленте $M$, нужно сдвинуть головку на ленте $M'$ на $k$ (число лент $M$) в одну из сторон. Значит, $1$ такт работы $M$ превратится в $k \cdot T(n)$ тактов на $M'$. А всё время работы $M$ превратится в $k \cdot (T(n))^2$ тактов на $M'$, то есть $O(T(n)) \rightarrow O((T(n)^2)$.
%\end{proof}

Теперь оценим замедление. Чтобы смоделировать один такт работы $M$ необходимо найти все индикаторы головок на ленте $M'$ и выполнить соответствующие изменения. В худшем случае на это потребуется полный проход всей ленты. Это означает, что один такт работы $M$ превратится в $k \cdot T(n)$ тактов на $M'$. А всё время работы $M$ превратится в $k \cdot (T(n))^2$ тактов на $M'$, то есть $O(T(n)) \rightarrow O((T(n)^2)$.
\end{proof}

\begin{corollary} Машину $M$ с ленточным алфавитом $\Gamma$ и с $k$ бесконечными в обе стороны лентами можно смоделировать на машине $M'$ с 1-ой бесконечной вправо лентой с замедлением в квадратичное число раз.
\end{corollary}

\begin{lemmanote} (б/д, не часть билета, но bear in mind) Машину $M$ с ленточным алфавитом $\Gamma$ и с $k$ бесконечными вправо лентами можно смоделировать на машине $M'$ с \emph{двумя} бесконечными вправо лентами с замедлением: $O(T(n)) \rightarrow O((T(n)\log(T(n))$.

Поэтому именно переход от двух лент к одной требует квадрат времени (на это задача 1).
\end{lemmanote}


\begin{lemma} Машину $M$ с ленточным алфавитом $\Gamma$ с $k$ бесконечными вправо лентами можно смоделировать на машине $M'$ с ленточным алфавитом $\{0, 1, \#\}$ и $k$ бесконечными вправо лентами с замедлением в const раз.
\end{lemma}

\begin{proof} Старую $\Gamma$ закодируем символами $\{0, 1\}$ одной длиной $l$, а разделителем будет служить $\#$. Тогда множество состояний увеличится в $l$ раз, и нужно будет считывать по $l$ символов за раз. Значит, замедление будет в $l$ раз, то есть в const раз.
\end{proof}

\begin{corollary} Машину $M$ с ленточным алфавитом $\Gamma$ и с $k$ бесконечными в обе стороны лентами можно смоделировать на машине $M'$ с ленточным алфавитом $\{0, 1, \#\}$ и с 1-ой бесконечной вправо лентой с полиномиальным замедлением.
\end{corollary}

%Идея в том, чтобы символы $\Gamma$ закодировать $\log|\Gamma|$ (одним числом) символами $\{0, 1\}$. Чтобы смоделировать шаг работы $M$, машина $M'$ должна: 

%\begin{enumerate}
%    \item за $\log|\Gamma|$ шагов прочитать с каждой из $k$ лент $\log|\Gamma|$ символов, кодирующих символ $\Gamma$.
%    \item используя состояния, запомнить прочитанные символы.
%    \item используя функцию перехода $M$, вычислить новые символы, которые пишет $M$, и новое состояние.
%    \item запомнить это всё в состоянии.
%    \item за $\log|\Gamma|$ шагов написать закодированные новые символы.
%\end{enumerate}

%Из конечности множества $Q$ состояний $M$ и алфавита $\Gamma$ следует, что множество состояний $Q'$ состояний $M'$ также конечно (Арора и Барак в своей книге почему-то оценивают сверху числом $10|Q||\Gamma|^kk$, непонятно, откуда 10). Также можно показать (можно, но не будем), что на любом входе $x \in \{0, 1\}^n$ если машина $M$ делает $T(n)$ шагов, то $M'$ делает не больше $4\log|\Gamma| T(n)$ шагов (то есть $\mathcal{O}(T(n))$ шагов)
